\subsubsection{代表性 OD 对路径形态比较}

为了更直观地解释四类算法的路径差异，本文选取 \verb|IN1 -> P5| 作为路径图案例。该 OD 对需要从入库区进入货架中部，路径需要绕开货架 A/F，能够展示栅格折线、可视图短路和转角代价评价之间的区别。另选 \verb|P8 -> OUT2| 作为表格补充，用于说明接近直线通行时栅格算法仍可能产生额外转向。

\begin{figure}[!htbp]
\centering
\includegraphics[width=0.86\linewidth]{p1_path_comparison_IN1_P5.png}
\caption{\texttt{IN1 -> P5} 的四类 P1 路径形态对比}
\label{fig:p1-compare-in1-p5}
\end{figure}

图 \ref{fig:p1-compare-in1-p5} 中，Grid A* 路径沿栅格方向逐步上行并右移，距离为 6.843，转角数为 5；VG 直接利用可视边绕过货架，距离降至 6.492，转角数为 2。AVG 与 DAVG 在这个 OD 上选择了和 VG 相同的几何折线，但由于 AVG/DAVG 的 \verb|total_cost| 中包含转角惩罚，其总成本为 6.729，高于 VG 的 6.492。这说明同一条几何路径在不同算法的评价口径下可能有不同总成本，P2 在读入路径源时需要明确使用的是 \verb|travel_time| 还是 \verb|total_cost|。

作为补充，\verb|P8 -> OUT2| 在连续空间中几乎可以直线到达，因此 VG、AVG、DAVG 均输出距离 3.583、转角数 0 的直线路径。Grid A* 由于受到 $0.1$m 栅格和 8 邻域运动方向限制，输出距离 3.863、转角数 18 的折线。这个案例说明，即使起终点之间没有复杂绕障碍需求，栅格路径也可能因为离散化而产生额外转向；连续可视图算法可以显著改善路径几何质量。

\begin{table}[H]
\centering
\small
\begin{tabular}{p{0.15\linewidth}p{0.13\linewidth}rrrr}
\toprule
OD 对 & 算法 & 距离 & 转角数 & 动态代价 & 总成本 \\
\midrule
\verb|IN1->P5| & Grid A* & 6.843 & 5 & 0.000 & 6.843 \\
\verb|IN1->P5| & VG & 6.492 & 2 & 0.000 & 6.492 \\
\verb|IN1->P5| & AVG & 6.492 & 2 & 0.000 & 6.729 \\
\verb|IN1->P5| & DAVG & 6.492 & 2 & 0.000 & 6.729 \\
\verb|P8->OUT2| & Grid A* & 3.863 & 18 & 0.000 & 3.863 \\
\verb|P8->OUT2| & VG & 3.583 & 0 & 0.000 & 3.583 \\
\verb|P8->OUT2| & AVG & 3.583 & 0 & 0.000 & 3.583 \\
\verb|P8->OUT2| & DAVG & 3.583 & 0 & 0.000 & 3.583 \\
\bottomrule
\end{tabular}
\caption{代表性 OD 对的路径指标对比}
\label{tab:p1-pair-metrics}
\end{table}

表 \ref{tab:p1-pair-metrics} 与图 \ref{fig:p1-compare-in1-p5} 相互印证：A* 的优势在于稳健可达和离散化处理，VG 的优势在于几何距离，AVG/DAVG 的优势在于用附加代价表达转向和风险偏好。对于上层调度而言，这些算法没有绝对优劣，关键是它们提供了不同路径风格。
